\section{Experiments}\label{sec:experiments}

This section presents numerical illustrations produced by the scripts in \texttt{experiments/ntkvariance/}. All scripts are headless (matplotlib ``Agg'' backend) and save figures next to the script; the figures below were copied into \texttt{refs/colt2026/figures/}. Setup: finite-width NTK scripts use a \textbf{3-layer (2 ReLU)} RF-LR architecture with hidden width 12000--20000 so that finite-width variance is small; deterministic proxy scripts iterate the recursion to large \(k\) (e.g.\ \(k_{\max}=4000\) or \(L=200\)). Bottleneck rank \(r\) is swept as indicated; \(n\) is the number of inputs, \(d\) the input dimension. Scripts in the first subsection are \textbf{deterministic}; the rest are \textbf{finite-width Monte Carlo} (they illustrate fluctuation scales and do not constitute proofs).

\subsection{Deterministic mean-field / proxy depth scaling (Theorem~\ref{thm:rflr-depth-scaling-mean}, Corollary~\ref{cor:conditioning-equicorrelated-highdim})}

\textbf{What is computed.} The deterministic ReLU EOC correlation recursion \(\rho_k=\varrho(\rho_{k-1})\) and the proxy NTK recursion \(\Theta^{(1)}(\rho_1)=1+s(\rho_1)\), \(\Theta^{(k)}(\rho_1)=1+\frac{1}{r}\Theta^{(k-1)}(\rho_1)\dot{s}(\rho_k)+\frac{1}{r}s(\rho_k)\) for \(k\ge 2\). Diagnostics: \(w_k/k\), \(k|\Theta^{(k)}-\Theta_\star(r)|\), \(rk(\Theta_{\mathrm{diag}}^{(k)}-\Theta_{\mathrm{off}}^{(k)})\). For the equicorrelated model, explicit eigenvalues \(\lambda_1=\Theta^{(L)}(1)+(n-1)\Theta^{(L)}(\rho_0)\), \(\lambda_\perp=\Theta^{(L)}(1)-\Theta^{(L)}(\rho_0)\) (e.g.\ \(n=64\), \(\rho_0=0\)).

\begin{figure}[htbp]
  \centering
  \includegraphics[width=0.71\textwidth]{confirm_thm15}
  \caption{Three panels (Theorem~\ref{thm:rflr-depth-scaling-mean}): (1) \(w_k/k\) vs \(k\) (stabilizes); (2) \(k|\Theta^{(k)}(\rho_1)-\Theta_\star(r)|\) vs \(k\) (bounded); (3) \(rk(\Theta_{\mathrm{diag}}^{(k)}-\Theta_{\mathrm{off}}^{(k)})\) vs \(k\) (stabilizes). One curve per rank \(r\in\{5,10,20,50\}\).}
  \label{fig:confirm_thm15}
\end{figure}

\begin{figure}[htbp]
  \centering
  \includegraphics[width=0.71\textwidth]{confirm_prop17}
  \caption{Equicorrelated spectrum (Corollary~\ref{cor:conditioning-equicorrelated-highdim}): (1) spike saturation \(L|\lambda_1/n-\Theta_\star(r)|/n\) vs \(L\) (log--log); (2) gap scaling \(rL\lambda_\perp\) vs \(L\) (log--log). One curve per \(r\).}
  \label{fig:confirm_prop17}
\end{figure}

\begin{figure}[htbp]
  \centering
  \includegraphics[width=0.64\textwidth]{decay_vs_depth}
  \caption{Product \(\prod_{k=\ell+1}^L \dot{\Sigma}^{(k)}/r\) vs \(j=L-\ell\) (semilogy), compared to the bound \((c_0/r)^j\). Exponential-in-depth suppression (Eq.~\ref{eq:ntk_explicit_form}).}
  \label{fig:decay_vs_depth}
\end{figure}

\begin{figure}[htbp]
  \centering
  \includegraphics[width=0.56\textwidth]{rho_vs_depth}
  \caption{\(\rho_k\) vs \(k\). Correlation aligns to 1 along depth (ReLU EOC path). From \texttt{scalingdepth.py} (\(L=200\), \(c_0=1/2\)).}
  \label{fig:rho_vs_depth}
\end{figure}

\begin{figure}[htbp]
  \centering
  \includegraphics[width=0.56\textwidth]{one_minus_rho_vs_depth_loglog}
  \caption{\(1-\rho_k\) vs \(k\) (log--log) with \(k^{-2}\) reference. Consistent with Theorem~\ref{thm:rflr-depth-scaling-mean}: \(1-\rho_k=O(k^{-2})\).}
  \label{fig:one_minus_rho_vs_depth_loglog}
\end{figure}

\subsection{Finite-width Monte Carlo: variance of a single NTK entry}

\textbf{What is computed.} Monte Carlo over random initializations of a 3-layer (2-ReLU) RF-LR network: \(\mathrm{Var}[K(-1,1)]\) vs \(r\) (log--log). Two inputs \(x=-1\), \(x'=1\); width 16000; 100 trials. \textbf{Paper link:} rank-driven concentration (Section~\ref{sec:no-rkhs-advantage}, Appendix concentration bounds).

\begin{figure}[htbp]
  \centering
  \includegraphics[width=0.6\textwidth]{variance_vs_r}
  \caption{\(\mathrm{Var}[K(-1,1)]\) vs \(r\) with \(\propto 1/r\) baseline. Entry-wise variance decreases with \(r\) (rank-driven concentration).}
  \label{fig:variance_vs_r}
\end{figure}

\subsection{Finite-width Monte Carlo: distribution of smallest centered eigenvalue}

\textbf{What is computed.} For a fixed random dataset on the unit sphere (\(n=64\), \(d=16\)), build the empirical Gram \(K\), center \(K_c=HKH\), and record \(\lambda_{\min}^+(K_c)\) (smallest strictly positive eigenvalue; centering forces one zero eigenvalue). Monte Carlo over 50 initializations; \(r\in\{5,10,20,50,100,200,500,1000\}\). Illustrates the remark that the smallest eigenvalue can be driven close to 0 by finite-width noise (Appendix~\ref{rmk:lambda-min-negative-appendix}).

\begin{figure}[htbp]
  \centering
  \includegraphics[width=0.68\textwidth]{min_eig_distribution}
  \caption{Histograms of \(\lambda_{\min}^+(HKH)\) over random initializations, for several \(r\). Distribution can be driven close to zero (matrix stays PSD).}
  \label{fig:min_eig_distribution}
\end{figure}

\begin{figure}[htbp]
  \centering
  \includegraphics[width=0.6\textwidth]{min_eig_mean_vs_r}
  \caption{Mean of \(\lambda_{\min}^+(HKH)\) vs \(r\) (log scale). Smallest positive eigenvalue increases or stabilizes with \(r\) (better conditioning).}
  \label{fig:min_eig_mean_vs_r}
\end{figure}

\subsection{Condition number: proxy vs empirical (equicorrelated and high-dim spherical)}

\textbf{What is computed.} \textbf{Equicorrelated:} Proxy Gram has \(\kappa_\perp=1\) on \(\mathbf{1}^\perp\) (Corollary~\ref{cor:conditioning-equicorrelated-highdim}). Script computes condition number of the \emph{centered} empirical Gram vs \(r\), mean and std over 40 initializations; \(n=32\), \(d=64\), \(\rho_0=0\). Validates concentration around 1 (Theorem~\ref{thm:equicorrelated-op-bound}). \textbf{High-dim spherical:} Same plot for i.i.d.\ uniform points on \(S^{d-1}\) with \(d=256,128\); \(\max_{i\neq j}|\rho_{ij}|=O(1/\sqrt{d})\) w.h.p., so data are approximately equicorrelated.

\begin{figure}[htbp]
  \centering
  \includegraphics[width=0.64\textwidth]{condition_number_proxy_vs_empirical}
  \caption{Condition number on \(\mathbf{1}^\perp\) vs \(r\), equicorrelated data. Proxy fixed at 1; empirical mean \(\pm\) std concentrates as \(r\) grows.}
  \label{fig:condition_number_proxy_vs_empirical}
\end{figure}

\begin{figure}[htbp]
  \centering
  \includegraphics[width=0.64\textwidth]{condition_number_highdim_spherical}
  \caption{Same \(\kappa\) vs \(r\) for i.i.d.\ uniform on \(S^{d-1}\), \(d=256\) and \(d=128\). \(\kappa_\perp\approx 1\), empirical \(\kappa\) concentrates (Corollary~\ref{cor:conditioning-equicorrelated-highdim}).}
  \label{fig:condition_number_highdim_spherical}
\end{figure}

\subsection{Condition number: non-equicorrelated data (Proposition~\ref{prop:condition-number-depth})}

\textbf{What is computed.} Condition number vs \(r\) for a \textbf{clustered} design (e.g.\ 4 clusters) with varying \(\rho_{ij}\). \(n=48\), \(d=64\), 30 trials; \(r\) swept to \(10^4\)--\(10^5\). Illustrates that \(\kappa\) need not approach 1 and the lower bound \(\kappa\ge\Omega(r\cdot L)\) can be large.

\begin{figure}[htbp]
  \centering
  \includegraphics[width=0.64\textwidth]{condition_number_non_equicorrelated}
  \caption{\(\kappa\) vs \(r\) for clustered (non-equicorrelated) data. Contrast with equicorrelated: \(\kappa\) can remain large (Proposition~\ref{prop:condition-number-depth}).}
  \label{fig:condition_number_non_equicorrelated}
\end{figure}

\subsection{Kernel regression risk and analytic Fisher/Kibble constants}

\textbf{Kernel regression:} Fix a target (linear or low-degree polynomial on the sphere); fit kernel ridge regression with the empirical RF-LR NTK Gram; plot test MSE vs \(r\). 64 train, 256 test, \(d=32\), 20 trials. Validates that as \(r\) grows and the Gram concentrates, risk improves. \textbf{Fisher/Kibble:} Closed-form curves (no sampling): Fisher density \(p(\hat{\rho}\mid \rho,r)\) (\({}_2F_1\)), Kibble joint density (modified Bessel \(I_\nu\)). \textbf{\(I(r)\) and spectral decay:} \(I(r)\) and \(\sqrt{r}\,I(r)\) vs \(r\); \(|2c_1(r)|\) (Puiseux spectral-decay amplitude). Paper: Remark~\ref{rmk:fisher-kibble-law}, Section~\ref{appendix:I-r-scaling}.

\IfFileExists{figures/kernel_regression_risk_vs_r.png}{%
\begin{figure}[htbp]
  \centering
  \includegraphics[width=0.6\textwidth]{kernel_regression_risk_vs_r}
  \caption{Test MSE of kernel ridge regression vs \(r\). Performance improves as the Gram concentrates toward the proxy.}
  \label{fig:kernel_regression_risk_vs_r}
\end{figure}%
}{}

\IfFileExists{figures/fisher_density_curves.png}{%
\begin{figure}[htbp]
  \centering
  \includegraphics[width=0.56\textwidth]{fisher_density_curves}
  \caption{Fisher density \(p(\hat{\rho}\mid \rho,r)\) (closed form). Analytic correlation distribution used in concentration (Remark~\ref{rmk:fisher-kibble-law}).}
  \label{fig:fisher_density_curves}
\end{figure}%
}{}

\IfFileExists{figures/kibble_density_contours.png}{%
\begin{figure}[htbp]
  \centering
  \includegraphics[width=0.56\textwidth]{kibble_density_contours}
  \caption{Kibble joint density contours (closed form). Radial part of Fisher--Kibble decoupling.}
  \label{fig:kibble_density_contours}
\end{figure}%
}{}

\IfFileExists{figures/I_of_r_scaling.png}{%
\begin{figure}[htbp]
  \centering
  \includegraphics[width=0.34\textwidth]{I_of_r_scaling}
  \hfill
  \includegraphics[width=0.34\textwidth]{sqrt_r_I_of_r}
  \caption{Left: \(I(r)\) vs \(r\). Right: \(\sqrt{r}\,I(r)\) vs \(r\). \(I(r)\sim C/\sqrt{r}\); \(\sqrt{r}\,I(r)\) stabilizes (Section~\ref{appendix:I-r-scaling}).}
  \label{fig:I_of_r_scaling}
\end{figure}%
}{}

\IfFileExists{figures/spectral_decay_prefactor_vs_r.png}{%
\begin{figure}[htbp]
  \centering
  \includegraphics[width=0.6\textwidth]{spectral_decay_prefactor_vs_r}
  \caption{\(|2c_1(r)|\) (spectral-decay amplitude) vs \(r\). Puiseux/RKHS \(r\)-dependent prefactor.}
  \label{fig:spectral_decay_prefactor_vs_r}
\end{figure}%
}{}
